% Appendix A — Alternative model comparison
% Updated 2026-05-19: main text adopts the per-subject selected loss
% (gamma + L_RDM; Methods, parameter selection). Retains LOCO-rho-argmax fits and the
% Machado/R+C cone-shift baselines as historical reference; R+C
% decomposition retained only as an exploratory descriptive companion
% (near-degenerate loss; no etiological claim).

\section*{Appendix A: Alternative Model Comparison}
\label{app:altmodels}

The main text adopts a 2-component cortical model (Section~\ref{sec:results:twocomp}).
Here we report two cone-shift baselines as reference comparisons:

\begin{itemize}
  \item \textbf{Machado (1-DOF)} --- a single retinal cone-spectral shift $\Delta\lambda$ \citeNP{machado2009}.
  \item \textbf{Retinal + Cortical (R+C, 2-DOF)} --- retinal shift $\Delta\lambda$ plus a cortical compensation gain $g$ \citeNP{tregillus2021}.
\end{itemize}

\subsection*{A.1 Fit results (LOCO Spearman $\rho$ argmax, historical reference)}

Under the LOCO-$\rho$ argmax criterion used in earlier drafts, the Machado model fit Sub-09 at $\Delta\lambda = 13.5$~nm with $\rho_\text{V4} = 0.76$ and Sub-08 at $\rho_\text{V4} = 0.62$ (lower fit quality); the R+C model did not improve on either single-component model in either subject.
These fits are reported as historical reference for the main-text 2-component fit (Section~\ref{sec:results:twocomp}); the LOCO-$\rho$ argmax criterion used here need not agree with the per-subject loss selected in the main text (\S\ref{sec:methods:selection}).

\paragraph{R+C decomposition as a diagnostic.}
Although the R+C 2-stage model is not adopted as the filter form (see \nameref{sec:results:filter}), the R+C parameter fit at hV4 is reported here as an exploratory per-subject diagnostic.
For sub-08, the R+C decomposition returns $\Delta\lambda = 2.5$~nm with $g = -2.25$; for sub-09 it returns $\Delta\lambda = 19.5$~nm with $g = -1.10$.
These values are reported for completeness only and carry no etiological interpretation.
The R+C neural loss is weakly identified (near-degenerate) along the $\Delta\lambda$--$g$ trade-off, and the $g$ estimate in particular is unstable across fitting variants, reversing sign under alternative criteria (the historical LOCO-$\rho$ argmax fit returns $g = +2.25$ for sub-08, versus $g = -2.25$ here; Supplementary Table~\ref{tab:S1}); the $(\Delta\lambda, g)$ pair therefore does not constitute a reliable point estimate and supports no retinal-versus-cortical attribution.
We retain it only as a descriptive companion to the 2-component fit, not as an independent test.

\subsection*{A.2 Inverse-mapping (pre-image) behavior}

The angular form of the 2-component model is bijective by construction: a stimulus-space pre-image $\theta + \delta\theta$ exists for any target hue.
The 1-DOF Machado fit for sub-09, by contrast, collapses three hues (green $135^\circ$, cyan $180^\circ$, blue $225^\circ$) onto a single pre-image angle ($\sim 127^\circ$), which mathematically prevents exact inversion at those positions.
This arc-compression failure is a structural property of the 1-DOF retinal-shift model and is independent of the fitting criterion; it disqualifies the Machado baseline for the stimulus-space filter task of Section~\ref{sec:results:filter}.

\subsection*{A.3 R+C model}

The R+C model is retained here as a 2-DOF baseline.
Because the cortical gain term $g$ acts as a uniform amplitude scaling without generating the angular distortion captured by the 2-component model's $\beta_c \cos(h - \theta_{\rm conf})$ term, R+C does not match the observed hV4 LOCO vulnerability pattern beyond what the retinal $\Delta\lambda$ term alone provides; this conclusion is consistent across loss criteria.
